\documentclass[10pt,a4paper]{amsart}
\usepackage[margin=19mm]{geometry}
\usepackage{amsmath,amssymb,mathtools}
\usepackage[scr=rsfs,cal=euler]{mathalfa}
\usepackage{xfrac}
\usepackage[unicode,hidelinks,bookmarks=false]{hyperref}
\newcommand{\ie}{i.e.\ }
\newcommand{\cf}{cf.\ }
\newcommand{\resp}{respectively\ }
\newcommand{\Subsection}{\subsection}
\providecommand{\nobreakdash}{\nobreak\hbox{-}}
\setlength{\emergencystretch}{2em}
\multlinegap0pt
\usepackage{amssymb}
\newtheorem{lemma}{Lemma}
\newcommand{\abs}[1]{\left|#1\right|} % Absolute value
\newcommand{\bN}{\mathbb{N}} % Natural numbers
\newcommand{\bR}{\mathbb{R}} % Reals
\newcommand{\cC}{\mathcal{C}} % Differentiable
\newcommand{\cS}{\mathcal{S}} % Schwartz space
\newcommand{\dv}{V} % Horizontal vector field
\newcommand{\dw}{W} % Vertical vector field
\newcommand{\dz}{Z} % Central vector field
\newcommand{\norm}[1]{\left\|#1\right\|} % Norm
\title{Anisotropic spaces and nil-automorphisms}
\author{Oliver Butterley, Minsung Kim}
\date{Source: arXiv:2308.06630v3 (CC BY 4.0)}
\begin{document}
\maketitle

\noindent\emph{Source and licence.} Anisotropic spaces and nil-automorphisms. Version arXiv:2308.06630v3 (CC BY 4.0), distributed by arXiv under the Creative Commons Attribution 4.0 International licence, http://creativecommons.org/licenses/by/4.0/, as recorded in the arXiv licence metadata for this exact version and shown on the abstract page https://arxiv.org/abs/2308.06630v3. Full licence notice: \texttt{licenses/arxiv-2308.06630-CC-BY-4.0.txt}.\par\smallskip

\noindent\textit{Editorial scope.} Counted unit: the article's lemma lem:slide (source0.tex lines 434--442). The article's complete original proof of it is delivered with it (source0.tex lines 444--505, one \texttt{proof} environment of 62 lines). No mathematics is written, repaired or re-derived: the statement and the proof are the article's own lines, in the article's own notation, and no retained line is substituted or re-flowed.  Retained uncounted as the source-local context the proof needs: lines 411--418.  Nothing was omitted from any retained span, so the package carries no omission record.  No retained line contains a citation, so the wrapper carries no bibliography and no bibliography entry.  No retained line contains a figure environment, an image-inclusion command or drawing code, so no image is delivered and the package declares no asset dependency.  Editorial formatting only, in addition: an AMS \texttt{amsart} preamble that restates the article's own declarations and macros used by the retained lines, namely the environments \texttt{lemma} on separate counters, together with the standard packages those lines need.  The retained environments print in this document as Lemma 1 in order of appearance, whereas the article prints the same items with the numbering of its own full document.  The complete native source of the article as distributed in the arXiv e-print for this exact version is bundled unchanged as \texttt{sources/145567/source0.tex}.
\begin{lemma}
    \label{lem:norm-scale}
    Given \(q\in \bN\), there exists \(C>0\) such that, for all \(\eta \in \cC^{\infty}(\bR)\) supported in some interval \(A \subset \bR\) of length \(\abs{A} \geq 2\delta\) and for all \(h\in \cC^{\infty}_{N}\),
    \[
        \abs{\int \eta(t) \cdot  h\circ \varphi_{t}^{\dw}(m) \ dt}
        \leq C \abs{A} \norm{\eta}_{\cC^q} \abs{h}_{0,q}.
    \]
\end{lemma}

\begin{lemma}
    \label{lem:slide}
    Suppose that \(m \in M\), \(q\in\bN\), \(0<\epsilon<\delta\), and \(\eta \in \cS(I_{\delta})\), \(\norm{\eta}_{\cC^q} \leq 1\).
    Let \(\tilde{m} = \varphi_c^\dz \circ \varphi_b^\dw \circ \varphi_a^\dv (m)\) and \(\tilde{\eta}(t) = e^{- 2\pi i NK (c + at)} \eta(t)\).
    Then there exists \(C>0\) such that for all \(h \in  \cC^\infty_N\) and any \(a,b,c \in \bR\) with \(\abs{a}, \abs{b} \leq \epsilon\), %the following holds.
    \[
        \abs{\ell_{\eta,m}(h) - \ell_{\tilde{\eta},\tilde{m}}(h)} \leq C \epsilon \norm{h}_{1,q-1}.
    \]
\end{lemma}

\begin{proof}
    Rearranging the integral defining the linear functional,
    \begin{equation}
        \label{eq:slide1}
        \begin{aligned}
            \ell_{\tilde{\eta},\tilde{m}}(h)
             & = \int \tilde{\eta}(t) \cdot h \circ \varphi_{t+b}^{\dw} \circ \varphi_a^\dv \circ \varphi_c^\dz (m) \ dt                                     \\
             & = \int \tilde{\eta}(t) \cdot h \circ \varphi_{t}^{\dw} \circ \varphi_a^\dv \circ \varphi_c^\dz (m) \ dt                                       \\
             & \ \ + \int \left[\tilde{\eta}(t-b) - \tilde{\eta}(t)\right] \cdot h \circ \varphi_{t}^{\dw} \circ \varphi_a^\dv \circ \varphi_c^\dz (m) \ dt.
        \end{aligned}
    \end{equation}
    Continuing with the first part of the last term, recalling that \(\varphi_{t}^{\dw} \circ \varphi_a^\dv = \varphi_a^\dv  \circ \varphi_{t}^{\dw} \circ \varphi_{at}^\dz\),
    \begin{multline}
        \label{eq:slide2}
            \int \tilde{\eta}(t) \cdot h \circ \varphi_{t}^{\dw} \circ \varphi_a^\dv \circ \varphi_c^\dz (m) \ dt
              = \int \tilde{\eta}(t) \cdot h \circ \varphi_a^\dv  \circ \varphi_{t}^{\dw} \circ \varphi_{c+at}^\dz (m) \ dt                               \\
             = \int \tilde{\eta}(t) \cdot h \circ \varphi_{t}^{\dw} \circ \varphi_{c+at}^\dz (m) \ dt                                                    
              \ \ + \int \int_{0}^{a} \tilde{\eta}(t) \cdot \dv h \circ \varphi_{s}^{\dv} \circ \varphi_{t}^{\dw} \circ \varphi_{c+at}^\dz (m) \ ds \ dt.
    \end{multline}
    Since \(h \in \cC^{\infty}_{N}\) and hence \(h \circ \varphi_{t}^{\dw}  \in \cC^{\infty}_{N}\), by definition of \(\tilde{\eta}\),
    \begin{equation}
        \label{eq:slide3}
        \begin{aligned}
            \int \tilde{\eta}(t) \cdot h \circ \varphi_{t}^{\dw} \circ \varphi_{c+at}^\dz (m) \ dt
             & = \int e^{2\pi i NK (c + at)} \tilde{\eta}(t) \cdot h \circ \varphi_{t}^{\dw} (m) \ dt \\
             & = \int \eta(t) \cdot h \circ \varphi_{t}^{\dw} (m) \ dt                               \\
             & = \ell_{\eta,m}(h).
        \end{aligned}
    \end{equation}
    Combining the above equalities \eqref{eq:slide1}, \eqref{eq:slide2} and \eqref{eq:slide3}, we have shown that
    \begin{equation}
        \label{eq:slide4}
        \begin{aligned}
            \abs{\ell_{\eta,m}(h) - \ell_{\tilde{\eta},\tilde{m}}(h)}
             & \leq \abs{\int \left[\tilde{\eta}(t-b) - \tilde{\eta}(t)\right] \cdot h \circ \varphi_{t}^{\dw} \circ \varphi_a^\dv \circ \varphi_c^\dz (m) \ dt} \\
             & \ \ + \abs{\int \int_{0}^{a} \tilde{\eta}(t) \cdot \dv h \circ \varphi_{s}^{\dv} \circ \varphi_{t}^{\dw} \circ \varphi_{c+at}^\dz (m) \ ds \ dt}.
        \end{aligned}
    \end{equation}
    For convenience when bounding the first of these two terms, let \(\zeta(t) = \tilde{\eta}(t) - \tilde{\eta}(t-b)\).
    Observe that the support of \(\zeta\) is contained within an interval of length \(2(\delta + \epsilon) \leq 4\delta\).
    Moreover, for each $n\in \bN$,  
    \[\zeta^{(n)}(t) = \tilde{\eta}^{(n)}(t) - \tilde{\eta}^{(n)}(t-b) = \int_{t-b}^{t} \tilde{\eta}^{(n+1)}(s) \ ds,\] and so \(\norm{\zeta}_{\cC^{q-1}} \leq b \norm{\tilde{\eta}}_{\cC^{q}}\).
    By assumption \(\norm{\eta}_{\cC^q} \leq 1\) and so there is an upper bound for \(\norm{\tilde{\eta}}_{\cC^{q}}\), uniform in \(\eta\), depending only on \(q\).
    Consequently, using also Lemma~\ref{lem:norm-scale}, there exists \(C_1>0\) such that,
    \begin{equation}
        \label{eq:slide5}
        \abs{\int \left[\tilde{\eta}(t-b) - \tilde{\eta}(t)\right] \cdot h \circ \varphi_{t}^{\dw} \circ \varphi_a^\dv \circ \varphi_c^\dz (m) \ dt} 
        \leq \epsilon C_1 \norm{h}_{0,q-1}.
    \end{equation}
    For the final term  \eqref{eq:slide4} which remains to estimate, 
    \begin{multline*}
        \int \left( \int_{0}^{a} \tilde{\eta}(t) \cdot \dv h \circ \varphi_{s}^{\dv} \circ \varphi_{t}^{\dw} \circ \varphi_{c+at}^\dz (m) \ ds \right) \ dt \\
        = \int_{0}^{a} \left( \int \tilde{\eta}(t) e^{2\pi i (c + at + st)N} \cdot \dv h  \circ \varphi_{t}^{\dw} \circ \varphi_{s}^{\dv} (m) \ dt \right) \ ds \\
        \leq a \norm{h}_{1,{q}} \norm{\smash{t \mapsto  \tilde{\eta}(t) e^{2\pi i (c + at + st)N} }}_{\cC^{q}}.
    \end{multline*}
    Consequently, there exists a constant \(C_2>0\), depending only on \(N\) and \(q\) such that this term is bounded above by
    \begin{equation}
        \label{eq:slide6}
        \epsilon C_2 \norm{h}_{1,q}.
    \end{equation}
    Combining the two estimates we have obtained, \eqref{eq:slide5} and \eqref{eq:slide6}, whilst noting that \(\norm{h}_{0,q-1} \leq \norm{h}_{1,q-1}\) and \(\norm{h}_{1,q} \leq \norm{h}_{1,q-1}\), completes the proof of the lemma.
\end{proof}

\end{document}
